\subsection{Loop Fayet-Iliopoulos terms in $T^2/Z_2$ models: instability and moduli stabilization --- arXiv:2003.03512}

\textbf{Authors:} Hiroyuki Abe, Tatsuo Kobayashi, Shohei Uemura, Junji Yamamoto

\paragraph{主な主張}
$T^2/Z_2$ 上の六次元超対称 Abelian gauge 理論で、局在 FI 項への 1-loop 補正と bulk zero mode の局在化が複素構造に依存し、その自己整合的な相殺条件によってモジュラスを固定し得ることを示した \cite{Abe:2020vmv}。

\paragraph{新規性}
傾いたトーラスを含む幾何へ解析を拡張し、真空の安定性を brane 電荷配置と複素構造の選択に結び付けた。

\paragraph{位置付け}
orbifold の局在量子効果から複素構造固定化を実現する機構であり、モジュラーポテンシャルを仮定する模型と比較できる。

\paragraph{コメント（読書メモ）}
$\tau=\omega$ に対応するモジュラス固定化の例として注目している。原典の表記では $\tau=e^{i\pi/3}=\omega+1$（$\omega=e^{2\pi i/3}$）であり、modular 変換で同値な点を表す。
